The equivalence of (108.11) and (108.12) may be checked by inserting definition (108.2) into (108.12) and using the orthogonality properties of the $3j$-symbols.

The high symmetry of the $9j$-symbol follows directly from definition (108.11) and the symmetry properties of the $3j$-symbols. Interchanging any two rows or any two columns multiplies the $9j$-symbol by $(-1)^{\sum j}$. It is also unchanged by transposition, that is, by interchanging rows and columns.

Scalars of still higher order depend on still more parameters $j$. Their number must clearly be a multiple of three; these are the $3nj$-symbols. We shall not consider their properties here. We mention only that for every $n>3$ there is more than one distinct type of $3nj$-symbol, and these types cannot be reduced to one another. Thus there are two distinct types of $12j$-symbols\footnote{For a fuller account of the theory of $9j$-symbols and of the properties of $3nj$-symbols, see the book by Edmonds cited on p.~272, and the books: A.~P.~Yutsis, I.~B.~Levinson, V.~V.~Vanagas. \emph{Mathematical Apparatus of the Theory of Angular Momentum}. --- Vilnius, 1960; D.~A.~Varshalovich, A.~N.~Moskalev, V.~K.~Khersonskii. \emph{Quantum Theory of Angular Momentum}. --- Moscow: Nauka, 1975.}.
